\section{Runtime Environment}
The production system is implemented in Haskell using GHC 9.4. Concurrency is built on GHC's lightweight threads (one OS thread per core, $\sim$60{,}000 green threads at peak) and STM for shared state. Redis access uses the \texttt{hedis} client; PostgreSQL access uses \texttt{postgresql-simple}.

\section{Thread-Level Cache}
\begin{lstlisting}[language=Haskell, caption=Thread-level cache (request-scoped)]
newtype ThreadCache = ThreadCache (IORef (HashMap CacheKey CacheVal))

newThreadCache :: IO ThreadCache
newThreadCache = ThreadCache <$> newIORef HM.empty

threadLookup :: ThreadCache -> CacheKey -> IO (Maybe CacheVal)
threadLookup (ThreadCache ref) k = HM.lookup k <$> readIORef ref

threadInsert :: ThreadCache -> CacheKey -> CacheVal -> IO ()
threadInsert (ThreadCache ref) k v = modifyIORef' ref (HM.insert k v)
\end{lstlisting}

\section{Server-Local Cache (TVar + LRU + TTL)}
\begin{lstlisting}[language=Haskell, caption=Server-local cache (node-scoped)]
data ServerCache = ServerCache
  { scStore   :: TVar (HashMap CacheKey (POSIXTime, CacheVal))
  , scMaxSize :: Int
  }

serverLookup :: ServerCache -> CacheKey -> IO (Maybe CacheVal)
serverLookup (ServerCache st _) k = do
  now <- getPOSIXTime
  atomically $ do
    m <- readTVar st
    case HM.lookup k m of
      Just (expiry, v) | expiry > now -> pure (Just v)
                      | otherwise     -> do
                          writeTVar st (HM.delete k m)
                          pure Nothing
      Nothing -> pure Nothing
\end{lstlisting}

\section{Integrated Read Path with Single-Flight}
\begin{lstlisting}[language=Haskell, caption=Two-tier read path with single-flight coalescing]
fetch :: Ctx -> CacheKey -> IO CacheVal
fetch ctx k = do
  -- Tier 1
  threadLookup (ctxThread ctx) k >>= \case
    Just v  -> pure v
    Nothing -> do
      -- Tier 2 + single-flight on miss
      v <- singleFlight (ctxFlight ctx) k $
             serverLookup (ctxServer ctx) k >>= \case
               Just v -> pure v
               Nothing -> do
                 v <- redisGet (ctxRedis ctx) k
                 serverInsert (ctxServer ctx) k v
                 pure v
      threadInsert (ctxThread ctx) k v
      pure v
\end{lstlisting}

\section{Authoritative Decrement Path}
\begin{lstlisting}[language=Haskell, caption=Authoritative inventory decrement using Lua]
decrementStock :: Connection -> ByteString -> Int -> IO (Either StockErr Int)
decrementStock conn sku qty = do
  let script = "local a = tonumber(redis.call('GET', KEYS[1]) or '0') \
               \ local w = tonumber(ARGV[1]) \
               \ if a >= w then redis.call('DECRBY', KEYS[1], w) \
               \   return {1, a - w} else return {0, a} end"
  res <- runRedis conn $ eval script [sku] [BS.pack (show qty)]
  case res of
    Right [Integer 1, Integer remaining] -> do
      publish conn ("invalidate:" <> sku) ""
      pure (Right (fromIntegral remaining))
    Right [Integer 0, Integer current]   ->
      pure (Left (Insufficient (fromIntegral current)))
    _ -> pure (Left RedisError)
\end{lstlisting}

\section{Cache Invalidation via Pub/Sub}
On successful authoritative writes the writer publishes an invalidation event on a Redis Pub/Sub channel. Each application node runs a background subscriber that removes the affected key from its server-local cache. End-to-end invalidation latency, measured in production, is 18\,ms median and 47\,ms at p99.

\section{Cross-Language Notes}
The same architecture maps cleanly onto other runtimes:
\begin{itemize}[leftmargin=*]
    \item \textbf{Java}: Caffeine for Tier 2; \texttt{ThreadLocal} for Tier 1; \texttt{singleflight} library or Caffeine's \texttt{AsyncCache} for coalescing; Redisson for Redlock.
    \item \textbf{Go}: \texttt{sync.Map} or \texttt{ristretto} for Tier 2; \texttt{context.Context} carrying a per-request map for Tier 1; \texttt{golang.org/x/sync/singleflight}; \texttt{redsync} for Redlock.
    \item \textbf{Rust}: \texttt{DashMap} or \texttt{moka} for Tier 2; per-task task-local for Tier 1; \texttt{singleflight} crate; \texttt{rslock} for Redlock.
    \item \textbf{Node.js}: \texttt{node-cache} or \texttt{lru-cache} for Tier 2; \texttt{AsyncLocalStorage} for Tier 1; \texttt{p-memoize} or custom promise map for coalescing; \texttt{redlock} package.
\end{itemize}

% =============================================================================
